\letterheader{\CandidateCityDE}

\CompanyName\\
\HiringTeamDE\\
\CompanyCityDE

\vspace{10pt}
\ApplicationDateDE

\vspace{10pt}
\textbf{Bewerbung als \JobTitleDE}

\vspace{8pt}
Sehr geehrtes Recruiting-Team,

% CUSTOMIZE: Jeden Absatz auf die Ausschreibung zuschneiden; maximal eine Seite.
hiermit bewerbe ich mich als \JobTitleDE. Mein Profil verbindet Dairy Technology, Dairy Science und Agrartechnik mit mehr als vier Jahren Erfahrung in der angewandten Nutztierforschung. Derzeit schließe ich am Leibniz-Institut für Agrartechnik und Bioökonomie (ATB) und an der Freien Universität Berlin meine Promotion zur Entwicklung skalierbarer, kostengünstiger Gassensorsysteme für das Emissionsmonitoring in der Nutztierhaltung ab.

An der Position reizt mich besonders die Verbindung von wissenschaftlichen Studien, Produktinnovation, Praxisbetrieben und Ergebnisvermittlung. Am ATB plane und realisiere ich Feldforschung in Milchviehställen, koordiniere Versuchsarbeiten mit Wissenschaftlern und technischen Mitarbeitenden und werte komplexe Langzeitdatensätze mit R aus. Ich war an vier Forschungsprojekten beteiligt, habe Praktikanten betreut und Publikationen sowie Konferenzbeiträge erstellt.

Mein M.Sc. in Dairy Science vermittelte mir eine fundierte Grundlage in Wiederkäuerernährung, Herdengesundheit, Precision Livestock Farming, Futterqualität, Milchwirtschaft und biometrischer Versuchsplanung. In meiner Forschung verknüpfe ich Emissions- und Stallklimadaten mit Herdenmanagementdaten. Dadurch kann ich technische Messungen im Zusammenhang mit Tierleistung und betrieblicher Praxis beurteilen.

Meine Promotion ist bewusst anwendungsorientiert: Ich entwickle und validiere kosteneffiziente Sensortechnologien gegenüber Referenzmessverfahren mit Blick auf ihre praktische Skalierbarkeit. In die Projekte von \CompanyName{} bringe ich sorgfältige Studienplanung, statistische Auswertung, Praxiserfahrung in Milchviehbetrieben und verständliche Wissenschaftskommunikation ein. Zugleich möchte ich meine Kenntnisse im Futtermittelrecht, in der Marktanalyse und in produktspezifischen Entwicklungsprozessen gezielt vertiefen.

Ich besitze eine Niederlassungserlaubnis und kann ohne Sponsoring in Deutschland arbeiten. Meine Deutsch- und Englischkenntnisse entsprechen dem Niveau B2; Reisebereitschaft ist vorhanden. Prof. Dr. Thomas Amon und Dr.-Ing. David Janke vom ATB stehen als Referenzen zur Verfügung. Über die Gelegenheit, Ihnen meinen möglichen Beitrag zu \CompanyName{} persönlich vorzustellen, freue ich mich.

Mit freundlichen Grüßen

\vspace{14pt}
\CandidateName

